% !PRINTED 502
\weberpage{502}
%<<CONT

3. \gesperrt{The absolute value of a product of two or
more rational functionals is equal to the product of the
absolute values of the factors.}

If $\omega$ is a functional of the field $\Omega$, then $N(\omega)$ is a rational
functional, and the \gesperrt{norm of a product or of a quotient
of two functionals is equal to the product or the quotient
of the norms of the parts.}

4. \gesperrt{We call the absolute value of the rational
functional $N(\omega)$ the absolute norm $N_a(\omega)$ of
$\omega$ and set}
\weq{3}{N(\omega) = N_a(\omega)\,E(\omega).}
$N_a(\omega)$ is always a positive rational number, and $E(\omega)$
is the quotient of two primitive integral functions. From 3.,
if $\alpha, \beta$ are two functionals in $\Omega$, there follows the formula
\weq{4}{N_a(\alpha\beta) = N_a(\alpha)\,N_a(\beta),}
or the theorem:

5. \gesperrt{The absolute norm of a product of functionals
in $\Omega$ is equal to the product of the absolute norms of the
factors.}

If we replace all the coefficients of a functional in $\Omega$ by
the corresponding numbers of a field $\Omega_1$ conjugate to $\Omega$,
we obtain a \gesperrt{conjugate functional}. The entire
functional field $\overline{\Omega}$ thereby passes over into a
\gesperrt{conjugate functional field} $\overline{\Omega}_1$. Since every equation
between functionals of the field $\Omega$ is at bottom only the
combination of a series of equations between numbers of the
field $\Omega$, we have the theorem (Bd. I, \S.\,141):

6. \gesperrt{Every equation between functionals of the field
$\overline{\Omega}$ remains correct if in place of all the functionals the
corresponding elements of a conjugate field $\overline{\Omega}_1$
are put.}

If $t$ denotes a variable not occurring in $\omega$, then
the function $N(t-\omega)$ has the form
\weq{5}{\Phi(t) = t^n + A_1 t^{n-1} + A_2 t^{n-2} + \cdots + A_n,}
in which the coefficients $A_i$ are rational functionals, and this
function $\Phi(t)$ vanishes when $\omega$ is put for $t$. There is,
however, not merely one such function $\Phi(t)$ which for $t = \omega$
%CONT>>
